% Verbatim round-6 source extracts for manual long-record migration.
% Not standalone; do not input blindly. See 02_long_record_migrations.md.

% ===== ORIGINAL sec:lcm; destination is specified in the migration plan =====
\section{The running least common multiple as a product of pure powers}\label{sec:lcm}

\begin{proposition}[the running least common multiple]\label{res:lcm}
Let $p,q,r$ be pairwise distinct primes and $x\ge1$.  Then
$\Lprefix(x)=\hgt(x)$.
\end{proposition}

\begin{proof}
For divisibility in one direction, every smooth $n\le x$ has exponents bounded
by the corresponding integer logarithms, so $n\mid\hgt(x)$, and hence
$\Lprefix(x)\mid\hgt(x)$.  For the other, the three pure powers
$p^{\lfloor\log_p x\rfloor}$, $q^{\lfloor\log_q x\rfloor}$ and
$r^{\lfloor\log_r x\rfloor}$ are themselves smooth numbers not exceeding $x$,
so each divides $\Lprefix(x)$.  Distinct primes have coprime powers, so their
product divides $\Lprefix(x)$ as well.  The two divisibilities give equality.
\end{proof}

\begin{proposition}\label{res:cube}
For every $x\ge1$, $\hgt(x)\le x^{3}$.
\end{proposition}

\begin{proof}
Each of the three factors is a power of its base not exceeding $x$.
\end{proof}

The opposite estimate follows from
$x/b<b^{\lfloor\log_b x\rfloor}$ for each $b\in\{p,q,r\}$:
\[
 \frac{x^3}{pqr}<\hgt(x)\le x^3.
\]
In particular the reciprocal kernel is summable over the exponent lattice.


% ===== ORIGINAL sec:cells; destination is specified in the migration plan =====
\section{Constancy on logarithmic cells and the jump points}\label{sec:cells}

By Proposition~\ref{res:lcm} the running value depends on $x$ only through the
three integer logarithms.  It is therefore constant wherever none of them
changes and moves only where one of them does, and the next definition names
the sets on which they are all constant, so that the next two theorems can say
where the value stands still and by what factor it moves.

Say that $x$ and $y$ lie in the same \emph{logarithmic cell} when
$\lfloor\log_b x\rfloor=\lfloor\log_b y\rfloor$ for each of $b=p,q,r$: the
\lword{Erdos269/ThreePrimeRunningLcm.lean}{163}{SameThreePrimeLogCell}{cell
relation}.

\begin{theorem}[constancy on cells]\label{res:cell}
If $x,y\ge1$ lie in the same logarithmic cell, then
$\Lprefix(x)=\Lprefix(y)$.  The same holds for the kernel at two smooth points
whose values lie in one cell.
\end{theorem}

\begin{proof}
Immediate from Proposition~\ref{res:lcm}, since $\hgt$ depends on $x$ only
through the three integer logarithms.
\end{proof}

\noindent Formalised as the
\lword{Erdos269/ThreePrimeRunningLcm.lean}{179}{smoothPrefixLcm_eq_of_sameLogCell}{cell
constancy of the running value}, the
\lword{Erdos269/ThreePrimeRunningLcm.lean}{169}{threePrimeHeight_eq_of_sameLogCell}{cell
constancy of the height}, and the
\lword{Erdos269/ThreePrimeRunningLcm.lean}{191}{threePrimeKernelQ_eq_of_sameLogCell}{cell
constancy of the kernel}.  The running value therefore changes only when one of
the three logarithms changes, and the next theorem says by exactly how much.

\begin{theorem}[single-coordinate jump ratios]\label{res:jump}
Let $x,y\ge1$.  If $\lfloor\log_p y\rfloor=\lfloor\log_p x\rfloor+1$ while the
other two logarithms agree, then $\Lprefix(y)=p\,\Lprefix(x)$; and similarly
with $q$ or $r$ in place of $p$.
\end{theorem}

\begin{proof}
By Proposition~\ref{res:lcm} both sides are heights, and one factor of the height
gains one in its exponent while the others are unchanged.
\end{proof}

\noindent Formalised as the
\lword{Erdos269/ThreePrimeRunningLcm.lean}{326}{smoothPrefixLcm_firstLogStep}{first},
\lword{Erdos269/ThreePrimeRunningLcm.lean}{339}{smoothPrefixLcm_secondLogStep}{second},
and
\lword{Erdos269/ThreePrimeRunningLcm.lean}{352}{smoothPrefixLcm_thirdLogStep}{third}
coordinate steps, over the corresponding statements for the height alone, which
need no primality: the
\lword{Erdos269/ThreePrimeRunningLcm.lean}{294}{threePrimeHeight_firstLogStep}{first
height step}, the
\lword{Erdos269/ThreePrimeRunningLcm.lean}{305}{threePrimeHeight_secondLogStep}{second},
and the
\lword{Erdos269/ThreePrimeRunningLcm.lean}{316}{threePrimeHeight_thirdLogStep}{third}.

The jump points are therefore the pure prime powers, and they do not collide
across channels.

\begin{theorem}[jump count]\label{res:count}
Let $n\ge0$.  The set of the first $n$ positive powers of $p$, of $q$, and of
$r$ has exactly $3n$ elements, and adjoining the common origin $1$ gives
exactly $3n+1$.
\end{theorem}

\begin{proof}
The count is routine.  Within one channel the powers $b,b^{2},\ldots,b^{n}$ are distinct because
$b\ge2$.  Across two channels a common value would be a positive power of two
distinct primes, which unique factorisation forbids.  Finally $1$ is not a
positive power of any prime.
\end{proof}

\noindent Formalised as the
\lword{Erdos269/ThreePrimeRunningLcm.lean}{249}{threePrimePositiveJumpSet_card}{positive
jump count} and the
\lword{Erdos269/ThreePrimeRunningLcm.lean}{278}{threePrimeJumpSetWithOrigin_card}{jump
count with the origin}, over the
\lword{Erdos269/ThreePrimeRunningLcm.lean}{208}{positivePrimePowers_card}{channel
cardinality}, the
\lword{Erdos269/ThreePrimeRunningLcm.lean}{229}{positivePrimePowers_disjoint}{channel
disjointness}, and the
\lword{Erdos269/ThreePrimeRunningLcm.lean}{219}{one_not_mem_positivePrimePowers}{exclusion
of the origin}; the channels themselves are the
\lword{Erdos269/ThreePrimeRunningLcm.lean}{204}{positivePrimePowers}{positive
power sets}.  The exponent $0$ is omitted from each channel because it is the
shared initial value, which is why the origin is counted once rather than three
times.

Two channels never meet, so at each positive pure power exactly one of the three
logarithms advances, and by Theorem~\ref{res:jump} the running value is
multiplied there by the prime of that channel.  Reading those multipliers in
increasing order of the pure powers gives the \emph{jump word} of $\Lprefix$:
one letter from $\{p,q,r\}$ for each positive pure power, recording which prime
the value is multiplied by at that point.  At $\{2,3,5\}$ the pure powers in
increasing order are $2,3,4,5,8,9,16,25,27,32,\ldots$, so the jump word begins
\[
 2,\;3,2,\;5,2,\;3,2,\;5,3,2,\;\ldots
\]
where the grouping is the one used in Section~\ref{sec:escape}: each group ends
at a power of two.


% ===== ORIGINAL sec:fibre; destination is specified in the migration plan =====
\section{The height-fibre normal form}\label{sec:fibre}

Grouping a lattice sum by the value of the running least common multiple turns
it into a sum over heights whose coefficients are multiplicities.  We prove
this exactly, on a finite rectangular box.

Write $\mathcal B(h_p,h_q,h_r)$ for the box of exponent triples with $i\le h_p$,
$j\le h_q$, $k\le h_r$, the
\lword{Erdos269/ThreePrimeRunningLcm.lean}{368}{smoothExponentBox}{exponent
box}, and, for a height $H$, write $F(H)$ for the set of points of the box whose
running height $\hgt(p^{i}q^{j}r^{k})$ equals $H$, the
\lword{Erdos269/ThreePrimeRunningLcm.lean}{377}{smoothHeightFiber}{height
fibre}.  Distinct lattice points can carry the same height, and the coefficients
of the normal form are exactly the sizes of these fibres.

\begin{theorem}[finite normal form]\label{res:fibre}
For every box $\mathcal B=\mathcal B(h_p,h_q,h_r)$,
\[
 \sum_{(i,j,k)\in\mathcal B}\Kernel(i,j,k)
 =\sum_{H}\frac{\#F(H)}{H},
\]
the outer sum ranging over the heights actually attained on $\mathcal B$.
\end{theorem}

\begin{proof}
The identity is a regrouping.  Partition $\mathcal B$ into the fibres of the
height map.  On $F(H)$ every summand is
$1/H$ by definition of the kernel, so the fibre contributes $\#F(H)/H$.
\end{proof}

\noindent Formalised as the
\lword{Erdos269/ThreePrimeRunningLcm.lean}{406}{finiteSmoothKernelSum_groupedByHeight}{height-fibre
normal form}, over the
\lword{Erdos269/ThreePrimeRunningLcm.lean}{384}{smoothHeightFiber_kernel_sum}{fibre
sum}, with the height of a lattice point given by the
\lword{Erdos269/ThreePrimeRunningLcm.lean}{373}{smoothPointHeight}{point
height}.

For a worked instance take $(p,q,r)=(2,3,5)$ and the box $\mathcal B(1,1,1)$,
whose eight points carry the smooth values $1,2,3,5,6,10,15,30$ and the heights
\[
\begin{array}{c|cccccccc}
p^{i}q^{j}r^{k}&1&2&3&5&6&10&15&30\\ \hline
\hgt&1&2&6&60&60&360&360&10800
\end{array}
\]
Six heights occur, two of them twice: the points $5$ and $6$ share the height
$60$, and $10$ and $15$ share the height $360$.  Theorem~\ref{res:fibre} here
reads
\[
 1+\tfrac12+\tfrac16+\tfrac1{60}+\tfrac1{60}+\tfrac1{360}+\tfrac1{360}
 +\tfrac1{10800}
 =1+\tfrac12+\tfrac16+\tfrac2{60}+\tfrac2{360}+\tfrac1{10800}
 =\tfrac{18421}{10800}.
\]
The two coefficients $2$ carry the whole content of the regrouping on this box;
where the heights are pairwise distinct the identity is a relabelling and
nothing more.

Together with Theorem~\ref{res:cell} this is the finite core of the ordered
prime-power jump expansion: the value is constant on cells, the cells are
indexed by heights, and the coefficient of a height is the number of lattice
points it collects.  The infinite shell identity is proved in Section~\ref{sec:actual-orbit}.  Theorem~\ref{res:fibre} is an identity between two finite sums,
and it is stated over the full box rather than the smooth prefix, so it is not
a statement about $\Lprefix$ at a cutoff.

The same module records a one-step map $\tau(b,d,s)=b(s-d)$, the
\lword{Erdos269/ThreePrimeRunningLcm.lean}{487}{tailStateStep}{variable-base
tail step}, with the rewriting
\lword{Erdos269/ThreePrimeRunningLcm.lean}{490}{tailStateStep_eq}{that names its
expanded form}.  No orbit of this map is analysed.


% ===== ORIGINAL sec:shell; destination is specified in the migration plan =====
\section{A quadratic bound for smooth exponent shells}\label{sec:shell}

A tail estimate needs to know how many lattice points can share a short
multiplicative interval.  Fix a box $\mathcal B(h_p,h_q,h_r)$ and an interval
$[\lo,\hi)$, and write $\mathcal S$ for the set of exponent triples of that box
whose smooth value lies in that interval, the
\lword{Erdos269/ThreePrimeRunningLcm.lean}{500}{smoothExponentShell}{smooth
exponent shell}.  Every bound below assumes the interval short in the
multiplicative sense, meaning that $\hi$ is at most a stated multiple of $\lo$.

Informally, the next lemma says that an interval whose right endpoint is at
most $b$ times its left endpoint contains at most one of the numbers $b^{a}w$.

\begin{lemma}[uniqueness in a short interval]\label{res:short}
Let $b\ge1$ and suppose $\hi\le b\,\lo$.  If $b^{a}w$ and $b^{a'}w$ both lie in
$[\lo,\hi)$, then $a=a'$.
\end{lemma}

\begin{proof}
If $a<a'$ then $\hi\le b\,\lo\le b\cdot b^{a}w=b^{a+1}w\le b^{a'}w<\hi$, which
is impossible; the case $a>a'$ is symmetric.
\end{proof}

\noindent Formalised as the
\lword{Erdos269/ThreePrimeRunningLcm.lean}{509}{exponent_unique_in_short_interval}{short-interval
uniqueness}.  The hypothesis is that the interval has multiplicative width at
most $b$.

\begin{proposition}\label{res:drop}
Let $\mathcal S$ be the shell of the box $\mathcal B(h_p,h_q,h_r)$ in
$[\lo,\hi)$.  If $\hi\le r\,\lo$ then $\#\mathcal S\le(h_p+1)(h_q+1)$; if
$\hi\le p\,\lo$ then $\#\mathcal S\le(h_q+1)(h_r+1)$.
\end{proposition}

\begin{proof}
Suppose $\hi\le r\,\lo$.  If two triples of $\mathcal S$ agree in their first
two coordinates, Lemma~\ref{res:short} applied with $b=r$ and $w=p^{i}q^{j}$
forces their third coordinates to agree as well.  So the projection forgetting
the third coordinate is injective on $\mathcal S$, and its image lies in a
rectangle with $(h_p+1)(h_q+1)$ points.  The case $\hi\le p\,\lo$ is the same
with the first coordinate projected away.
\end{proof}

\noindent Formalised as the
\lword{Erdos269/ThreePrimeRunningLcm.lean}{585}{smoothExponentShell_card_le_dropThird}{third-coordinate
projection} and the
\lword{Erdos269/ThreePrimeRunningLcm.lean}{543}{smoothExponentShell_card_le_dropFirst}{first-coordinate
projection}.

\begin{theorem}[quadratic multiplicity bound]\label{res:shell}
Let $\mathcal S$ be the shell of the box $\mathcal B(h_p,h_q,h_r)$ in
$[\lo,\hi)$.  Suppose $\hi\le r\,\lo$ and $h_p\le h_q\le h_r$ with
$h_p+h_q+h_r=j$.  Then
\[
 9\,\#\mathcal S\le(j+3)^{2}.
\]
\end{theorem}

\begin{proof}
By Proposition~\ref{res:drop}, $\#\mathcal S\le(h_p+1)(h_q+1)$; under the
sorting hypothesis the two surviving coordinates are the two smallest.  It
therefore
suffices to prove that $a\le b\le c$ with $a+b+c=j$ gives
$9(a+1)(b+1)\le(j+3)^{2}$.  From $a\le b\le c$ we get $a+2b\le j$, so it is
enough that $9(a+1)(b+1)\le(a+2b+3)^{2}$; writing $b=a+d$ with $d\ge0$, the
difference of the two sides is $d(3a+4d+3)\ge0$.
\end{proof}

\noindent Formalised as the
\lword{Erdos269/ThreePrimeRunningLcm.lean}{646}{smoothExponentShell_card_quadratic}{quadratic
shell bound}, over the
\lword{Erdos269/ThreePrimeRunningLcm.lean}{629}{sorted_pair_quadratic}{sorted
quadratic estimate}.  The constant $9$ is the square of the number of
generating primes and appears because the bound is the arithmetic--geometric
comparison for a sum of three sorted coordinates; sorting is a hypothesis, not
a normalisation, since the shell itself is not symmetric in the three bases.
The last inequality of the proof is an equality when $h_p=h_q=h_r$, both sides
then being $9(h_p+1)^{2}$, so no constant larger than $9$ survives that step.

The bound is uniform in $\lo$ and $\hi$ subject to the width condition, and it
is stated for the actual filtered shell rather than for a lattice model of it.
It is an input to a tail estimate and is not itself one: no series is bounded
here.  The estimate is elementary and uses no analytic input on the
distribution of smooth numbers, only the projection of
Proposition~\ref{res:drop}.


% ===== ORIGINAL sec:actual-orbit; destination is specified in the migration plan =====
\section{The actual shell orbit and its integral branch}\label{sec:actual-orbit}

We now pass from finite kernel structure to the literal infinite
$\{2,3,5\}$ shell sum.  Let $s_a$ be the total reciprocal running-height mass
of the smooth exponent triples in $[2^a,2^{a+1})$, put
\[
 T_a=\sum_{n\ge0}s_{a+n},\qquad
 X_a=\frac{\hgt(2^a)}2T_a,
\]
and let $b_a\in\{2,6,10,30\}$ and $m_a\in\N$ be respectively the ordered
block radix and ordered block digit.  These are the literal objects, not an
abstract bounded-radix model.

\begin{theorem}[actual shell-orbit dichotomy]\label{res:actual-orbit}
The series $\sum_{a\ge0}s_a$ is summable, and for every $a\ge0$,
\[
 X_{a+1}=b_aX_a-m_a.
\]
Moreover, either $X_a\in\Z$ for some $a$, or for every $a_0$ there is
$a\ge a_0$ such that
\[
 |X_a-z|\ge\frac1{31}\qquad\hbox{for every }z\in\Z.
\]
\end{theorem}

The recurrence follows by splitting the infinite tail after its first shell
and using the exact height and ordered-digit identities.  Summability follows
from the quadratic shell-count bound and geometric dyadic decay.  The last
assertion is the bounded-radix alternative applied with $2\le b_a\le30$.
It does not eliminate the first branch.  The same alternative applies to
every positive integer multiple $Y_a=BX_a$: the recurrence
$Y_{a+1}=b_aY_a-Bm_a$ keeps the radices in $\{2,6,10,30\}$, so either some
$BX_a$ is integral, and then every later term is, or
$\operatorname{dist}(BX_a,\mathbb Z)\ge 1/31$ cofinally.  A
distance-to-integer mandate that excludes the integral branch by hypothesis
therefore leaves the remaining throat untouched.

The jump word of $\{2,3,5\}$ cannot contain three consecutive factors of $2$:
between $2^e$ and $2^{e+2}$ the integer $v=\lfloor\log_3(2^e)\rfloor+1$
satisfies $2^e<3^v<2^{e+2}$, by the defining property of the integer
logarithm.  That unique-factorisation fact is the supplied-source three-power interval
lemma (\texttt{exists\_three\_power\_strictly\_between\_two\_powers} at
\path{ErdosProblems/Erdos269/JumpConstraintMajorant.lean}:28).  The geometric majorant $Q(n)=(n^2+8n+18)/9$ used only that
every multiplier is at least $2$.  Weighting the remaining product by the
stronger constraint that at least $\lfloor(k+1)/3\rfloor$ of the next $k$
letters are at least $3$, and taking the generating function
$A(z)=(1+z/2+z^2/6)/(1-z^3/12)$, yields the strictly smaller quadratic
\[
 \widetilde Q(n)=\frac{1210n^2+9130n+18847}{11979}.
\]
The moment identities and the comparison $Q-\widetilde Q>0$ are the
supplied-source moment form
(\texttt{carryMajorantQtilde\_eq\_moments} at
\path{ErdosProblems/Erdos269/JumpConstraintMajorant.lean}:102) and the
strict improvement (\texttt{carryMajorantQtilde\_lt\_Q} at line 119).  Assembling the same shell-multiplicity argument as
Theorem~\ref{res:shell} with these weights gives $X_a\le\widetilde Q(n_a)$,
where $n_a$ counts the positive $\{2,3,5\}$-powers not exceeding $2^a$.  An
exact truncated comparison, with remainder bounded by the already-recorded
checker width $a^2+6a+11$, is
\texttt{ErdosProblems/Erdos269/check\_jump\_constraint\_majorant.py}.  This
does not change the two directions of the window-escape criterion
separately.  Irrationality implies cofinal escape for every short bound
dominated by $c(B)(n+1)^2$, $\widetilde Q$ included.  The reverse
implication uses a \emph{valid} carry-dominating cap; the registered
Lean width $90B(a+1)^2$ is such a cap, and the actual if-and-only-if is
stated for that cap.  A generic upper-quadratic bound alone is not an
equivalence: the zero cap makes least-positive-residue escape automatic.

The dyadic clearing used below is one member of a symmetric family.

\begin{theorem}[prime-power boundary clearing]\label{res:prime-power-boundary-clearing}
Let $s\in\{2,3,5\}$, let $m\ge0$, and let $x$ be a positive
$\{2,3,5\}$-smooth integer with $x<s^m$.  Then
\[
 s\,\hgt(x)\mid\hgt(s^m).
\]
\end{theorem}

\begin{proof}
Write $x=2^i3^j5^k$.  The strict inequality forces the exponent in the
$s$-channel to be at most $m-1$; monotonicity of the other two integer
logarithms already places their exponents below those of $\hgt(s^m)$.
Divisibility is therefore coordinatewise.  The exact three-channel statement
is the supplied-source checked prime-power boundary clearing theorem
(\texttt{smoothHeight\_mul\_prime\_dvd\_boundaryHeight} at
\path{ErdosProblems/Erdos269/RationalLatticeReduction.lean}:51).
\end{proof}

The $s=2$ case supplies the normaliser for the shell orbit.  The $3$- and
$5$-power boundaries are additional exact clearing families, but this paper
does not yet couple them to their own normalised tail recurrences.  Write
$h_a=\hgt(2^a)/2$, which is an integer for $a\ge1$.

\begin{theorem}[all-scale rationality lattice]\label{res:all-scale-lattice}
Every finite shell window clears at its upper normaliser: for all
$u,v\ge0$ there is $m\in\N$ such that
\[
 h_{u+v}\sum_{i=0}^{v-1}s_{u+i}=m.
\]
If $q>0$ and $T_1=p/q$ for integers $p,q$, then
\[
 qX_a\in\Z\qquad(a\ge1).
\]
Consequently there are $i<j$ and $z\in\Z$ with
\[
 X_{1+j}-X_{1+i}=z.
\]
\end{theorem}

\begin{proof}
Every denominator in the finite window divides the height at its upper
endpoint, which proves the first assertion.  Splitting $T_1$ into that finite
window and the tail at $a$, multiplying by $q h_a$, and using the clearing
identity proves $qX_a\in\Z$.  Among $q+1$ such integers two have the same
residue modulo $q$; their corresponding states differ by an integer.
\end{proof}

Both conclusions are checked as supplied-source labels:
\texttt{qsmul\_normalizedTailState\_eq\_int\_of\_value\_eq\_rat}
(integrality of $qX_a$,
\path{ErdosProblems/Erdos269/RationalLatticeReduction.lean}:196) and
\texttt{exists\_normalizedTailState\_collision\_of\_value\_eq\_rat}
(the collision, line 296).

Thus pairwise incongruence of the actual $X_a$ modulo one would prove
irrationality.  The source does not establish it: rationality predicts a
collision rather than contradicting any known orbit theorem.

The integral branch in Theorem~\ref{res:actual-orbit} has a further exact
structure.

\begin{theorem}[pinning and rigidity of the integral branch]\label{res:pinning}
For every $a\ge0$,
\[
 X_a=\frac{m_a}{b_a}+\frac{X_{a+1}}{b_a},\qquad X_a>0.
\]
If $X_a$ is integral, then $X_n$ is integral for every $n\ge a$.

For every $m,K\ge0$, the genuine state has the exact finite-depth expansion
\[
 X_m=
 \sum_{i=0}^{K-1}
   \frac{m_{m+i}}{\prod_{j=0}^{i}b_{m+j}}
 +h_mT_{m+K}.                                      \tag{14}
\]

More generally, fix $A$, a positive width function $w$, and a real orbit
$(y_n)_{n\ge A}$ satisfying $y_{n+1}=b_ny_n-m_n$.  Suppose that for every
$n\ge A$, both $y_n$ and $X_n$ lie in
\[
 \left(\frac{m_n}{b_n},\frac{m_n}{b_n}+w(n)\right],
\]
and that for every $\varepsilon>0$ there is $k_0$ such that
$w(A+k)/2^k<\varepsilon$ for all $k\ge k_0$.  Then $y_A=X_A$.
\end{theorem}

\begin{proof}
The first identity is the recurrence solved for $X_a$; positivity follows
because every shell contains the dyadic point $2^a$.  Integer coefficients in
the forward recurrence give upward closure.  Iterating the pinning identity to
depth $K$ gives the finite mixed-radix digit sum plus
$X_{m+K}/\prod_{j<K}b_{m+j}$.  The height telescope
$\hgt(2^{m+K})=\hgt(2^m)\prod_{j<K}b_{m+j}$ and
$X_{m+K}=h_{m+K}T_{m+K}$ turn that remainder exactly into $h_mT_{m+K}$,
proving~\emph{(14)}.  This is the supplied-source checked finite-depth tail telescope
(\texttt{trueNormalizedState\_eq\_telescope} at
\path{ErdosProblems/Erdos269/IntegralBranchWidth.lean}:24); no limiting
argument is hidden.

For the final assertion, the
difference of the two orbits is multiplied by $b_n\ge2$ at each step, whereas
the common window bounds its absolute value at time $A+k$ by $w(A+k)$.
Hence
\[
 2^k|y_A-X_A|\le w(A+k).
\]
The stated decay forces $y_A=X_A$.
\end{proof}

This theorem turns indefinite survival of an integral seed inside the exact
windows into equality with the genuine infinite tail.  It does not prove that
no seed survives.  The four theorem families in this section are
proved above, with the linked Lean declarations providing exact mechanical
checks of their statements.


% ===== ORIGINAL sec:open; includes independent Fourier theorem =====
\section{Complements and further questions}\label{sec:open}

Theorems~\ref{res:two-prime-transcendence}
and~\ref{res:two-prime-repeated-transcendence} close both two-prime questions,
at the stronger level of transcendence.  The de-duplicated three-prime series
is not an open irrationality question: Erd\H{o}s recorded that result for any
finite set of primes in the letter of 1 January 1973
\cite[p.~335]{erdos1974letter}, without printing a proof.  The repeated
three-prime series lies outside the one-dimensional Hecke--Mahler reduction
and remains open.  At three primes the exact unresolved repeated-series
statement is best separated from the bridge, from the historical de-duplicated
assertion, and from the possible methods for proving the repeated sum.

\subsection{The actual block data and the remaining escape}

The checker already uses literal data, which we record so that the open
problems have no unspecified forcing word.  For $p\in\{2,3,5\}$ let $q,r$ be
the other two primes and put
\[
 A_p(e)=\#\{(i,j)\in\N^2:q^i r^j<p^e\},\qquad
 C_p(e)=\sum_{u=1}^{e}A_p(u).
\]
Let $I_a$ be the increasing list of internal jumps $(p,e)$ with
$p\in\{3,5\}$ and $2^a<p^e<2^{a+1}$.  For $(p,e)\in I_a$ let
\[
 \sigma_a(p,e)=\prod_{(q,f)\in I_a,\ p^e<q^f}q.
\]
Then, for $a\ge1$, the exact checker definitions are
\[
 \beta_a=2\prod_{(p,e)\in I_a}p,\qquad
 m_a^{235}=A_2(a+1)+
   \sum_{(p,e)\in I_a}(p-1)\sigma_a(p,e)
       \bigl(C_p(e)-C_2(a)\bigr).
\tag{9.1}\label{eq:actual-digit}
\]

The first formula is the four-letter radix of
Theorem~\ref{res:dyadic-alphabet}; the second is the integer implemented by the
pinned checker.  If
\[
 \nu_a=\#\{p^e:p\in\{2,3,5\},\ e\ge1,\ p^e<2^{a+1}\},
\]
its exact tested short bound is
\[
 K^{235}(B,a)=
 \left\lfloor\frac{B(\nu_a^2+10\nu_a+27)}9\right\rfloor.
\tag{9.2}\label{eq:actual-bound}
\]
The raw shell tail is the $T_a$ of Section~\ref{sec:actual-orbit}.  The
normalised state $X_a=\frac{\hgt(2^a)}2T_a$ has the mixed-radix expansion
\[
 X_a=\sum_{j\ge a}
   \frac{m_j^{235}}{\beta_a\beta_{a+1}\cdots\beta_j};
 \qquad X_{a+1}=\beta_a X_a-m_a^{235}.
\tag{9.3}\label{eq:actual-tail}
\]

The former open problem \texttt{prob:bridge269} is closed.  Its content is
Theorem~\ref{res:denominator-reduction}, with registered onset
$a_0=u+1+2v+3w$ and Lean width $90B(a+1)^2$.  The generic carry theorems of
Section~\ref{sec:escape} are now identified with the original series at that
onset.  The checker bound $K^{235}$ remains a valid finite majorant inside
the same quadratic family; it is not required for the registered equivalence.

\subsection{The intrinsic tail question}

\begin{problem}[exact nonintegrality of every reduced tail]\label{prob:tails269}
For every $B\ge1$ with $\gcd(B,30)=1$ and every $a\ge1$, prove
\[
 B X_a\notin\Z.
\tag{9.4}\label{eq:tail-nonintegrality}
\]
\end{problem}

This pointwise form is stronger-looking but cleaner than “cofinally
nonintegral”: if $B X_a$ is integral at one index, the recurrence
\[
 B X_{a+1}=\beta_a B X_a-Bm_a^{235}
\]
makes it integral at every later index.  Thus a direct solution of
Problem~\ref{prob:tails269}, joined to the proved bridge of
Theorem~\ref{res:denominator-reduction}, bypasses all residue-window
machinery.

The bounded-radix theorem gives a useful exact reduction.  Since
$2\le\beta_a\le30$, any real affine tail orbit either hits an integer or is,
cofinally often, at distance at least $1/31$ from every integer
(\lword{Erdos269/BoundedRadixTailEscape.lean}{89}{boundedRadix_zero_or_cofinal_far}{checked}).
It does not exclude the integral branch; Problem~\ref{prob:tails269} is exactly
what must do so for the actual orbit.

\subsection{A denominator-adaptive sufficient criterion}

For the literal pair $(m^{235},K^{235})$, retain the window definitions of
Section~\ref{sec:escape}:
\[
 W_{\ell,h}=\prod_{j=0}^{h-1}\beta_{\ell+j},\qquad
 F_{\ell,0}=0,\qquad
 F_{\ell,h+1}=\beta_{\ell+h}F_{\ell,h}+m_{\ell+h}^{235}.
\]

\begin{problem}[actual cofinal local-window escape]\label{prob:producer}
Prove the displayed quantifier order
\[
 \forall B\ge1\ (\gcd(B,30)=1),\ \forall a_0\ge1,\
 \exists\ell\ge a_0\ \exists h\ge1:\quad
 \operatorname{lpr}_{W_{\ell,h}}(-BF_{\ell,h})
   >K^{235}(B,\ell+h).
\tag{9.5}\label{eq:actual-escape}
\]
\end{problem}

Every $\beta_a$ is positive, so $W_{\ell,h}>0$ is automatic.  The cofinal
quantifier is present for a substantive reason: the proved bridge supplies
the reduced recurrence only after the denominator-dependent onset $a_D$,
and~\eqref{eq:actual-escape} then supplies a window beyond that onset.
Once such a window is chosen, the positive endpoint carry must equal the
canonical residue exactly
(\lword{Erdos269/RestrictedFloorSum.lean}{497}{leastPositiveResidue_windowForcing_eq_carry}{checked}),
yet it lies in the possible carry set
$\{1,\ldots,K^{235}(B,\ell+h)\}$; the strict inequality excludes that set.
This is a one-sided least-positive-residue statement, not two symmetric arcs
around zero.

Two exact countermodels rule out tempting shortcuts.  For $(W,F,B)=(6,4,1)$,
\[
 \operatorname{lpr}_{6}(-4)=2,
\]
so the canonical residue need not be coprime to $W$.  For
$(W,F,B)=(60,47,37)$,
\[
 \operatorname{lpr}_{60}(-37\cdot47)=1,
\]
so a fixed window has no denominator-independent positive residue lower
bound.  Therefore the window may genuinely depend on $B$; neither a fixed
finite computation nor a universal bounded-length guess addresses the
unbounded-denominator and cofinal-start quantifiers.  Equivalence of those
quantifiers with irrationality, in the absence of a rate, is not an
algorithmic impossibility: it names the same arithmetic in another
coordinate.  The finite checker can reject a proposed sufficient condition,
but supplies no cofinal statement.

\subsection{An exact torus--Fourier representation of the repeated value}

The repeated three-prime value itself, rather than only its radix word, has
an explicit analytic coordinate.  For $s=2^i3^j5^k$ put
$L=\log_2s$, $\theta=\log_2 3$, and $\phi=\log_2 5$.  Let
$f_p(x)=p^{\{x\}}$ be one-periodic, with value $1$ at its jump, and let
$F_{p,M}$ be its symmetric Fourier partial sum over $|m|\le M$.

\begin{theorem}[rectangular torus--Fourier identity]
\label{res:torus-fourier-value}
The absolutely convergent sums
\[
 A_M=\sum_{i,j,k\ge0}
 \frac{F_{2,M}(L)F_{3,M}(L/\theta)F_{5,M}(L/\phi)}
      {8^i27^j125^k}
\]
satisfy
\[
 \lim_{M\to\infty}A_M
 =\sum_{i,j,k\ge0}\frac1{\hgt(2^i3^j5^k)}+\frac{17}{2}.
 \tag{9.6}\label{eq:torus-fourier-value}
\]
For each finite $M$, expansion into Fourier modes is exact.  If
$\lambda=m_1+m_2/\theta+m_3/\phi$ and
$\widehat f_p(m)=(p-1)/(\log p-2\pi i m)$, the $(m_1,m_2,m_3)$ mode is
\[
 \frac{\widehat f_2(m_1)\widehat f_3(m_2)\widehat f_5(m_3)}
 {(1-e(\lambda)/8)(1-e(\lambda\theta)/27)
  (1-e(\lambda\phi)/125)},
 \qquad e(x)=e^{2\pi i x}.
\]
\end{theorem}

\begin{proof}[Proof spine]
Floor--fractional-part decomposition gives the pointwise identity
\[
 \frac1{\hgt(s)}=
 \frac{f_2(L)f_3(L/\theta)f_5(L/\phi)}{s^3}.
\]
For fixed $M$ the mode expansion is finite, and the three lattice sums are
ordinary geometric series.  Each $f_p$ has bounded variation.  The
Dirichlet--Jordan estimate therefore bounds all $F_{p,M}$ uniformly in $M$
and gives midpoint convergence at the jump.  The summable majorant
$C\,8^{-i}27^{-j}125^{-k}$ permits dominated convergence on the exponent
lattice.

The midpoint differs from the original torus weight only on the three pure
prime axes and at the origin.  The origin contributes $8$.  At $p^n$ the
excess is $(p-1)/(2\hgt(p^n))$, while the exact jump telescope
\[
 \sum_{p\in\{2,3,5\}}(p-1)
       \sum_{n\ge1}\frac1{\hgt(p^n)}=1
\]
makes the three axes contribute $1/2$.  Their total correction is therefore
$17/2$.
\end{proof}

This is an exact ordinary Fourier-analysis theorem, not a numerical fit.
It replaces an opaque scalar value by an explicit symmetric triple-mode
limit whose frequencies are linear forms in $1,\log_3 2,\log_5 2$.
The triple Fourier series is not absolutely summable, however, and the
identity supplies no Diophantine lower bound on a linear form in the value.
The three geometric denominators are bounded away from zero by $7/8$,
$26/27$ and $124/125$; their convergence requires no small-divisor estimate
on the phase $\lambda$.  That phase is a ratio of logarithms, not
automatically an integer linear form in $\log 2,\log 3,\log 5$.  The identity
proves neither irrationality nor transcendence and carries no priority claim.
Its authored proof is
\path{ErdosProblems/Erdos269/TensorDirichletJordanCompletion.md}; the finite
Fourier computation is regression evidence rather than proof authority.

\subsection{A two-dimensional analytic representation}

For the de-duplicated series define
\[
 \mathcal D_{2,3,5}=1+
 \sum_{t\in\{2^n,3^n,5^n:n\ge1\}}
 \frac{1}{
   2^{\lfloor\log_2t\rfloor}
   3^{\lfloor\log_3t\rfloor}
   5^{\lfloor\log_5t\rfloor}}.
\]
Irrationality of $\mathcal D_{2,3,5}$ is the historical de-duplicated
assertion of the 1973 letter, not an open problem of this note.  The problem
below asks only for a function-level representation or a proof that none
exists in a specified class.
The dyadic coding is the joint rotation word
\[
 \delta_{3,a}=\lfloor(a+1)\theta_3\rfloor-\lfloor a\theta_3\rfloor,
 \quad \theta_3=\frac{\log2}{\log3},\qquad
 \delta_{5,a}=\lfloor(a+1)\theta_5\rfloor-\lfloor a\theta_5\rfloor,
 \quad \theta_5=\frac{\log2}{\log5},
\]
with $\beta_a=2\,3^{\delta_{3,a}}5^{\delta_{5,a}}$.

\begin{problem}[function-faithful two-dimensional representation]
Express $\mathcal D_{2,3,5}$ as a nonconstant algebraic combination of values
of a specified two-dimensional Hecke--Mahler, cone-generating or multivariate
Mahler function and verify every hypothesis of a published value theorem; or
give a conditional theorem under an explicit logarithmic nondegeneracy
hypothesis; or prove that the literal series has no representation in the
specified finite-dimensional class.
\end{problem}

Pairwise irrationality of $\theta_3$ and $\theta_5$ is not silently promoted
to the orbit-closure or equidistribution hypothesis a two-dimensional theorem
may need.  The finite-observer formalisation isolates the precise
faithfulness requirement: equality in a finite observer must imply equality
after symbolic realisation, and a genuine finite-dimensional factorisation
forces the realised symbolic span to be finite-dimensional
(\lword{Erdos269/WeightedPhaseCarry.lean}{109}{carry_eq_residueDigit_add_coboundary}{residue--coboundary form},
\lword{Erdos269/WeightedPhaseCarry.lean}{293}{CartierRealisation}{symbolic realisation},
\lword{Erdos269/WeightedPhaseCarry.lean}{334}{finite_realisedSpan_of_factorisation}{checked}).
The same finite formulation keeps the carry residue and residue digit in their
declared intervals
(\lword{Erdos269/WeightedPhaseCarry.lean}{150}{carryResidue_mem_interval}{carry interval},
\lword{Erdos269/WeightedPhaseCarry.lean}{157}{residueDigit_mem_interval}{digit interval}).
No theorem here proves that the literal realised span is infinite, so a
general finite separable decomposition is neither assumed nor declared
excluded.

\subsection{Further structural constraints}

The radix alphabet itself extends without difficulty.  For ordered primes
$p_1<\cdots<p_s$, an interval $(p_1^a,p_1^{a+1})$ contains at most one power
from each other channel, because consecutive $p_i$-powers have ratio
$p_i>p_1$.  Hence its block radix belongs to the $2^{s-1}$-letter alphabet
\[
 \left\{p_1\prod_{i=2}^s p_i^{\varepsilon_i}:
   \varepsilon_i\in\{0,1\}\right\}.
\]
The nontrivial question is quantitative: prove effective recurrence or
discrepancy for the actual four-letter $\{2,6,10,30\}$ word, an asymptotic
formula with an error term for the restricted two-dimensional shell counts
that generate $m_a^{235}$, or the exact finite-separation rank of the literal
three-prime kernel under a specified family of shifts.

The three-channel rigidity and carry-lift extinction theorems already exclude
one proposed argument in four exact steps.  Under channel surjectivity, ordinary
block-nullity is equivalent to the perturbation being a coboundary of a
channel potential
(\lword{Erdos269/ThreeChannelBlockRigidity.lean}{59}{channelBlockNull_iff_channelPotential}{potential
classifier}).  Zero perturbations on genuine $2\to3$ and $2\to5$ transitions
then identify all three potential values and force the perturbation to vanish
at every index.  For an integral lift, that vanishing makes a nonzero initial
lift error grow by the exact product of the successive bases; bases at least
two make its absolute value at least $2^N$, contradicting even a single
index-$N$ bound strictly below $2^N$
(\lword{Erdos269/CarryLiftExtinction.lean}{178}{no_carryLift_of_errorBound_below_twoPow}{one-index
extinction}).  A single index therefore already suffices, and consequently no
uniform bound on the lift error can hold either: under the same hypotheses a
nonzero initial error is incompatible with any bound valid at every index
(\lword{Erdos269/CarryLiftExtinction.lean}{238}{no_bounded_carryLift_of_blockNull_twoAnchors}{uniform
extinction}).

A separate four-state calculation reaches the obstruction earlier: four real
states in $(0,1)$ with unit-accuracy integral lifts and the two anchor
equalities force the first complete $2$-block sum to be $1$, hence that block
cannot be null
(\lword{Erdos269/CarryLiftExtinction.lean}{289}{binaryLift_twoAnchors_force_firstBlock_sum_one}{first-block
sum},
\lword{Erdos269/CarryLiftExtinction.lean}{308}{no_unitAccurateLift_with_twoAnchors_and_firstTwoBlockNull}{four-state
obstruction}).

These lift obstructions assume the two anchor equalities and block-nullity.
The actual orbit and its ordered-power word have already been constructed;
the specified faithful lift and its additional hypotheses have not.  The affine alternative above may produce an
arbitrary integral state, not necessarily zero; its cofinal $1/31$ separation
is neither eventual nor positive-density and gives no unbounded distance.
The four-state calculation supplies no cofinal windows, and
\texttt{rational\_of\_scaledTail\_integer} classifies no denominators.  The
actual carry supplies a weighted block defect instead, so any successful
argument must use that weighted identity or construct a different faithful
lift.  None of these results is an irrationality theorem for \#269.

The bridge is proved.  The unresolved arithmetic is the nonintegrality
of $BX_a$ for every $a\ge1$ and every positive $B$ coprime to $30$.
